\clearpage
\oldpage[IV-2]{226}
\phantomsection
\addcontentsline{toc}{section}{Índice terminológico}
\section*{ÍNDICE TERMINOLÓGICO}

\begingroup
\footnotesize
\setlength{\parindent}{0pt}
\setlength{\parskip}{1.1pt}
\newcommand{\EGAIVtwoTerm}[2]{#1: #2.\par}

\EGAIVtwoTerm{Álgebra separable sobre un cuerpo}{4.6.2}
\EGAIVtwoTerm{Anillo de series formales restringidas sobre un cuerpo valuado completo}{7.4.8}
\EGAIVtwoTerm{Anillo excelente}{7.8.2}
\EGAIVtwoTerm{Anillo formalmente catenario}{7.1.9}
\EGAIVtwoTerm{Anillo formalmente equidimensional}{7.1.1}
\EGAIVtwoTerm{Anillo geométricamente normal, geométricamente reducido,
geométricamente regular, que posee la propiedad $(R_k)$ geométrica}{6.7.6}
\EGAIVtwoTerm{Anillo normal}{5.13.5}
\EGAIVtwoTerm{$\boldsymbol P$-anillo}{7.3.13}
\EGAIVtwoTerm{Anillo que posee la propiedad $(R_k)$}{5.8.2}
\EGAIVtwoTerm{Anillo estrictamente equidimensional}{7.2.1}
\EGAIVtwoTerm{Anillo estrictamente formalmente catenario}{7.2.6}
\EGAIVtwoTerm{Anillo universalmente catenario}{5.6.2}
\EGAIVtwoTerm{Antifiltro de subconjuntos de un conjunto}{5.10.8}

\medskip
\EGAIVtwoTerm{Codimensión de un subconjunto cualquiera de un preesquema}{5.1.3}
\EGAIVtwoTerm{Coprofundidad de un módulo en un punto, coprofundidad de un módulo,
coprofundidad de un preesquema}{5.7.1}
\EGAIVtwoTerm{Cuerpo de definición}{4.8.4}
\EGAIVtwoTerm{Cuerpo separablemente cerrado}{4.3.3}
\EGAIVtwoTerm{Ciclo primo asociado a un módulo, a un preesquema}{3.1.1}
\EGAIVtwoTerm{Ciclo primo asociado inmerso}{3.1.1}

\medskip
\EGAIVtwoTerm{Descomposición irredundante, descomposición irredundante reducida
de un módulo}{3.2.5}
\EGAIVtwoTerm{Descomposición primaria de $0$ en un módulo}{3.2.5}
\EGAIVtwoTerm{Dimensión de un módulo}{5.1.12}
\EGAIVtwoTerm{Dimensión de un preesquema sobre un cuerpo}{4.1.1}

\medskip
\EGAIVtwoTerm{Exponente de inseparabilidad de una extensión de un cuerpo}{4.7.4}
\EGAIVtwoTerm{Extensión primaria de un cuerpo}{4.3.1}

\medskip
\EGAIVtwoTerm{Clausura algebraica separable de un cuerpo en una extensión}{4.3.4}
\EGAIVtwoTerm{Fibra formal de un anillo en un punto de su espectro}{7.3.13}

\medskip
\EGAIVtwoTerm{Homeomorfismo universal}{2.4.2}
\EGAIVtwoTerm{Homomorfismo de módulos definido sobre un cuerpo}{4.8.4}

\medskip
\EGAIVtwoTerm{Lema de Hironaka}{5.12.8}
\EGAIVtwoTerm{Longitud geométrica de un módulo en un punto}{4.7.5}

\medskip
\EGAIVtwoTerm{Módulo íntegro sobre un anillo}{3.2.4}
\EGAIVtwoTerm{$A$-módulo que posee la propiedad $(S_k)$}{5.7.2}
\EGAIVtwoTerm{Módulo reducido sobre un anillo}{3.2.2}
\EGAIVtwoTerm{$\mathcal O_X$-módulo de Cohen-Macaulay en un punto,
$\mathcal O_X$-módulo de Cohen-Macaulay}{5.7.1}
\EGAIVtwoTerm{$\mathcal O_X$-módulo fielmente plano (con respecto a un morfismo)}{2.2.4}
\EGAIVtwoTerm{$\mathcal O_X$-módulo geométricamente íntegro sobre un cuerpo}{4.6.19}

\clearpage
\oldpage[IV-2]{227}

\EGAIVtwoTerm{$\mathcal O_X$-módulo geométricamente íntegro en cada punto sobre un cuerpo}{4.6.22}
\EGAIVtwoTerm{$\mathcal O_X$-módulo geométricamente reducido (o separable) sobre un cuerpo}{4.6.17}
\EGAIVtwoTerm{$\mathcal O_X$-módulo geométricamente reducido (o separable) sobre un cuerpo en un punto}{4.6.22}
\EGAIVtwoTerm{$\mathcal O_X$-módulo equidimensional en un punto}{5.1.12}
\EGAIVtwoTerm{$\mathcal O_X$-módulo íntegro, íntegro en un punto}{3.2.4}
\EGAIVtwoTerm{$\mathcal O_X$-módulo irredundante}{3.2.4}
\EGAIVtwoTerm{$\mathcal O_X$-módulo normalmente plano a lo largo de un subpreesquema}{6.10.1}
\EGAIVtwoTerm{$\mathcal O_X$-módulo que posee la propiedad $(S_k)$,
que posee la propiedad $(S_k)$ en un punto}{5.7.2}
\EGAIVtwoTerm{$\mathcal O_X$-módulo reducido, reducido en un punto}{3.2.2}
\EGAIVtwoTerm{$\mathcal O_X$-módulo $Z$-cerrado, $Z$-puro}{5.9.9}
\EGAIVtwoTerm{$\mathcal O_X$-módulo $Z$-cerrado, $Z$-puro en un punto}{5.9.9}
\EGAIVtwoTerm{Morfismo birracional de preesquemas reducidos}{6.15.4}
\EGAIVtwoTerm{Morfismo conexo}{4.5.5}
\EGAIVtwoTerm{Morfismo de Cohen-Macaulay en un punto, morfismo de Cohen-Macaulay}{6.8.1}
\EGAIVtwoTerm{Morfismo de coprofundidad $\leq n$ en un punto, morfismo de coprofundidad $\leq n$}{6.8.1}
\EGAIVtwoTerm{Morfismo definido sobre un cuerpo}{4.8.4}
\EGAIVtwoTerm{Morfismo irreducible}{4.5.5}
\EGAIVtwoTerm{Morfismo suave en un punto, morfismo suave}{6.8.1}
\EGAIVtwoTerm{Morfismo normal en un punto, morfismo normal}{6.8.1}
\EGAIVtwoTerm{$\boldsymbol P$-morfismo}{7.3.1}
\EGAIVtwoTerm{Morfismo que posee la propiedad $(R_k)$ en un punto,
morfismo que posee la propiedad $(R_k)$}{6.8.1}
\EGAIVtwoTerm{Morfismo que posee la propiedad $(S_k)$ en un punto,
morfismo que posee la propiedad $(S_k)$}{6.8.1}
\EGAIVtwoTerm{Morfismo cuasifielmente plano, morfismo cuasiplano}{2.3.3}
\EGAIVtwoTerm{Morfismo radicial en un punto}{6.15.3}
\EGAIVtwoTerm{Morfismo reducido en un punto, morfismo reducido}{6.8.1}
\EGAIVtwoTerm{Morfismo regular en un punto, morfismo regular}{6.8.1}
\EGAIVtwoTerm{Morfismo de desingularización}{7.9.1}
\EGAIVtwoTerm{Morfismo universalmente abierto, universalmente bicontinuo}{2.4.2}
\EGAIVtwoTerm{Multiplicidad radicial, multiplicidad separable, multiplicidad
total de una extensión de un cuerpo}{4.7.4}
\EGAIVtwoTerm{Multiplicidad radicial, multiplicidad separable, multiplicidad
total de un preesquema sobre un cuerpo}{4.7.4}
\EGAIVtwoTerm{Multiplicidad radicial de un punto para un $\mathcal O_X$-módulo
($X$ preesquema sobre un cuerpo)}{4.7.5}
\EGAIVtwoTerm{Multiplicidad radicial de un ciclo primo maximal asociado a un
$\mathcal O_X$-módulo ($X$ preesquema sobre un cuerpo)}{4.7.8}
\EGAIVtwoTerm{Multiplicidad total de un punto (de un ciclo primo maximal
asociado) para un $\mathcal O_X$-módulo ($X$ preesquema sobre un cuerpo)}{4.7.12}

\medskip
\EGAIVtwoTerm{Número geométrico de componentes conexas, de componentes
irreducibles, para un preesquema sobre un cuerpo}{4.5.2}

\medskip
\EGAIVtwoTerm{Punto asociado a un $\mathcal O_X$-módulo, a un preesquema}{3.1.1}
\EGAIVtwoTerm{Punto de Cohen-Macaulay de un preesquema}{5.7.1}
\EGAIVtwoTerm{Punto geométricamente unirrama, punto unirrama de un preesquema}{6.15.1}
\EGAIVtwoTerm{Preesquema conexo en codimensión $d\setmin 1$}{5.10.8}
\EGAIVtwoTerm{Preesquema de Cohen-Macaulay}{5.7.1}
\EGAIVtwoTerm{Preesquema excelente}{7.8.5}
\EGAIVtwoTerm{Preesquema (sobre un cuerpo) geométricamente conexo}{4.5.2}
\EGAIVtwoTerm{Preesquema (sobre un cuerpo) geométricamente íntegro}{4.6.2}
\EGAIVtwoTerm{Preesquema (sobre un cuerpo) geométricamente íntegro en cada punto}{4.6.9}
\EGAIVtwoTerm{Preesquema (sobre un cuerpo) geométricamente irreducible}{4.5.2}
\EGAIVtwoTerm{Preesquema (sobre un cuerpo) geométricamente normal en un punto,
geométricamente normal}{6.7.6}
\EGAIVtwoTerm{Preesquema (sobre un cuerpo) geométricamente reducido}{4.6.2}
\EGAIVtwoTerm{Preesquema (sobre un cuerpo) geométricamente reducido en un punto}{4.6.9}
\EGAIVtwoTerm{Preesquema (sobre un cuerpo) geométricamente regular en un punto,
geométricamente regular}{6.7.6}

\clearpage
\oldpage[IV-2]{228}

\EGAIVtwoTerm{Preesquema geométricamente unirrama en un punto,
geométricamente unirrama}{6.15.1}
\EGAIVtwoTerm{Preesquema íntegro en un punto}{3.2.4}
\EGAIVtwoTerm{Preesquema irredundante}{3.2.4}
\EGAIVtwoTerm{Preesquema que admite localmente una inmersión en un esquema regular}{5.11.1}
\EGAIVtwoTerm{Preesquema íntegro en cada punto}{4.6.9}
\EGAIVtwoTerm{Preesquema que posee la propiedad $(R_k)$, que posee la propiedad
$(R_k)$ en un punto}{5.8.2}
\EGAIVtwoTerm{Preesquema (sobre un cuerpo) que posee la propiedad $(R_k)$
geométrica en un punto, que posee la propiedad $(R_k)$ geométrica}{6.7.6}
\EGAIVtwoTerm{Preesquema que posee la propiedad $(S_k)$, que posee la propiedad
$(S_k)$ en un punto}{5.7.2}
\EGAIVtwoTerm{Preesquema regular en codimensión $k$}{5.8.2}
\EGAIVtwoTerm{Preesquema (sobre un cuerpo) separable}{4.6.2}
\EGAIVtwoTerm{Preesquema (sobre un cuerpo) separable en un punto}{4.6.9}
\EGAIVtwoTerm{Preesquema unirrama en un punto, preesquema unirrama}{6.15.1}
\EGAIVtwoTerm{Preesquema universalmente catenario}{5.6.3}
\EGAIVtwoTerm{Preesquema (sobre un cuerpo) universalmente reducido}{4.6.2}
\EGAIVtwoTerm{Profundidad de un $\mathcal O_X$-módulo en un punto}{5.7.1}
\EGAIVtwoTerm{Propiedad $\boldsymbol P$ del primer tipo (del segundo tipo) asociada a una
propiedad $\boldsymbol R$ de anillos locales noetherianos}{7.3.10}
\EGAIVtwoTerm{Propiedad $\boldsymbol P$ geométrica}{7.3.6}

\medskip
\EGAIVtwoTerm{Resolución de singularidades}{7.9.1}

\medskip
\EGAIVtwoTerm{Subconjunto (de un preesquema) definido sobre un cuerpo}{4.8.4}
\EGAIVtwoTerm{Sub-$\mathcal O_X$-módulo definido sobre un cuerpo}{4.8.4}
\EGAIVtwoTerm{Sub-$\mathcal O_X$-módulo primario de un $\mathcal O_X$-módulo}{3.2.4}
\EGAIVtwoTerm{Subpreesquema definido sobre un cuerpo}{4.8.4}
\EGAIVtwoTerm{Subpreesquema cerrado primario de un preesquema}{3.2.4}

\medskip
\EGAIVtwoTerm{$Z$-clausura de un $\mathcal O_X$-módulo}{5.9.11}

\endgroup

\par\smallskip
\noindent\textit{Nota de lectura.} Se conserva el orden de las entradas del
índice francés impreso; no se presenta como orden alfabético español.
Los números remiten a los apartados originales de EGA.
